% AQA Harvard Referencing Style Template
% For English Academic Writing
\documentclass[12pt, a4paper, twoside]{article}

% Essential Packages
\usepackage{geometry}
\usepackage{graphicx}
\usepackage{amsmath}
\usepackage{amsfonts}
\usepackage{amssymb}
\usepackage{hyperref}
\usepackage{harvard}
\usepackage{booktabs}
\usepackage{caption}
\usepackage{subcaption}
\usepackage{enumitem}
\usepackage{setspace}

% Page Settings
\geometry{
  left=2.5cm,
  right=2.5cm,
  top=2.5cm,
  bottom=2.5cm,
  headheight=15pt,
  headsep=1cm,
  footskip=1cm
}

% Line Spacing
\onehalfspacing

% Title Settings
\title{\textbf{Paper Title}}
\author{Author Name}
\date{\today}

\begin{document}

% Title Page
\maketitle

% Abstract
\begin{abstract}
  This paper introduces... (Abstract content)
\end{abstract}

% Keywords
\textbf{Keywords:} Keyword 1, Keyword 2, Keyword 3

% Table of Contents
\tableofcontents
\newpage

% Main Body
\section{Introduction}
Classical mechanics is deterministic in form, yet many deterministic systems are not predictably simple in practice. Small differences in initial conditions can produce strongly divergent trajectories in non-linear dynamics, so mathematically well-defined equations may still generate motion that appears irregular over finite observation windows. This contrast between formal structure and physical complexity motivates the present investigation.

The double pendulum is an appropriate system for this purpose. It consists of two point masses connected by two rigid, massless strings, with the upper mass attached to a pivot. In the small-angle limit, its behaviour can be analysed through linear normal-mode theory; outside that limit, non-linear coupling leads to richer responses. In this project, the focus is not chaos as an isolated topic, but the broader non-linear response of a \emph{periodically driven} double pendulum, including transitions between regular and irregular motion as control parameters vary.

The study addresses three linked questions. First, how accurately do small-angle analytical predictions describe measured or simulated frequencies in weak-amplitude motion? Second, how do forcing frequency and forcing amplitude influence the qualitative response of the driven system? Third, to what extent does including experimentally identified damping improve agreement between theoretical and numerical results? These questions establish a direct connection between analytical mechanics, numerical modelling, and empirical parameter estimation.

The methodology integrates theory, simulation, and experiment. A Lagrangian formulation is used to derive the equations of motion, and linearisation is applied to obtain the small-angle modal frequencies as reference quantities. Numerical integration is then performed in Python with a Runge-Kutta scheme, while parameter sweeps are carried out over selected ranges of initial conditions and driving parameters. In parallel, a single-pendulum experiment is used to estimate the local damping coefficient, which is subsequently inserted into the driven double-pendulum model. System response is assessed through time series, phase-space trajectories, and frequency-domain analysis, allowing quantitative validation of the small-angle approximation and critical evaluation of model limitations.

The remainder of this report is organised as follows. Section 2 reviews relevant literature on non-linear oscillators, driven pendulum dynamics, and deterministic sensitivity. Section 3 presents the theoretical framework, including Lagrangian derivation, linear modal analysis, and superposition. Section 4 describes numerical implementation and experimental damping identification. Section 5 reports simulation and validation results, and Section 6 concludes with limitations and future extensions.

\section{Literature Review}
  Literature review content...

\section{Theoretical Analysis}
To thoroughly investigate the dynamics of the double pendulum, a robust mathematical model must be established. This section derives the equations of motion using Lagrangian mechanics, applies small-angle approximations to find analytical solutions, and explores the physical implications of the system's normal modes.

\subsection{Kinematics and the Lagrangian}
Consider a double pendulum consisting of two point masses, $m_1$ and $m_2$, suspended by two massless rigid rods of lengths $l_1$ and $l_2$. Let $\theta_1$ and $\theta_2$ be the angles the rods make with the vertical downward direction. The Cartesian coordinates of the masses, taking the suspension point as the origin and the downward direction as the positive $y$-axis, are given by:
\begin{align*}
x_1 &= l_1 \sin \theta_1, \quad y_1 = l_1 \cos \theta_1, \\
x_2 &= l_1 \sin \theta_1 + l_2 \sin \theta_2, \quad y_2 = l_1 \cos \theta_1 + l_2 \cos \theta_2.
\end{align*}

The kinetic energy ($T$) of the system is the sum of the kinetic energies of both masses, and the potential energy ($V$) is determined by their vertical positions in a uniform gravitational field $g$:
\begin{align*}
T &= \frac{1}{2}m_1(\dot{x}_1^2 + \dot{y}_1^2) + \frac{1}{2}m_2(\dot{x}_2^2 + \dot{y}_2^2), \\
V &= -m_1 g y_1 - m_2 g y_2.
\end{align*}

By substituting the time derivatives of the coordinates into the kinetic energy equation, the Lagrangian of the system, defined as $L = T - V$, is constructed:
\begin{equation}
\begin{aligned}
L &= \frac{1}{2}(m_1 + m_2)l_1^2\dot{\theta}_1^2 + \frac{1}{2}m_2 l_2^2 \dot{\theta}_2^2 + m_2 l_1 l_2 \dot{\theta}_1 \dot{\theta}_2 \cos(\theta_1 - \theta_2) \\
  &\quad + (m_1 + m_2)g l_1 \cos\theta_1 + m_2 g l_2 \cos\theta_2.
\end{aligned}
\end{equation}

\subsection{Equations of Motion and Small-Angle Approximation}
The dynamics of the system are governed by the Euler-Lagrange equations:
\begin{align*}
\frac{d}{dt}\left(\frac{\partial L}{\partial \dot{\theta}_i}\right) - \frac{\partial L}{\partial \theta_i} &= 0, \quad \text{for } i = 1, 2.
\end{align*}

Evaluating these partial derivatives yields a set of highly non-linear, coupled second-order differential equations. For arbitrary angles, these equations exhibit chaotic behavior and can only be solved numerically. However, to understand the fundamental oscillatory properties, we analyze the system under the small-angle approximation, where $\theta_1, \theta_2 \ll 1$. Applying Taylor series expansions, we approximate $\sin\theta \approx \theta$ and $\cos\theta \approx 1 - \frac{\theta^2}{2}$, while ignoring higher-order terms such as $\dot{\theta}^2\theta$. The linearized equations of motion become:
\begin{equation}
\begin{aligned}
(m_1 + m_2)l_1 \ddot{\theta}_1 + m_2 l_2 \ddot{\theta}_2 + (m_1 + m_2)g \theta_1 &= 0, \\
m_2 l_1 \ddot{\theta}_1 + m_2 l_2 \ddot{\theta}_2 + m_2 g \theta_2 &= 0.
\end{aligned}
\end{equation}

\subsection{Derivation of Natural Frequencies}
To find the normal modes of the system, we seek oscillatory solutions of the form $\theta_1(t) = A_1 e^{i\omega t}$ and $\theta_2(t) = A_2 e^{i\omega t}$, where $\omega$ is the angular frequency and $A_1, A_2$ are the amplitudes. Substituting these into the linearized equations yields an eigenvalue problem:
\begin{equation}
\begin{pmatrix}
(m_1 + m_2)g - (m_1 + m_2)l_1\omega^2 & -m_2 l_2 \omega^2 \\
-m_2 l_1 \omega^2 & m_2 g - m_2 l_2 \omega^2
\end{pmatrix}
\begin{pmatrix}
A_1 \\ A_2
\end{pmatrix} = \begin{pmatrix} 0 \\ 0 \end{pmatrix}.
\end{equation}

For non-trivial solutions, the determinant of the coefficient matrix must be zero. Solving this characteristic equation gives the natural frequencies of the system. For simplicity, if we consider the special case where $m_1 = m_2 = m$ and $l_1 = l_2 = l$, the frequencies simplify to:
\begin{equation}
\omega_{\pm}^2 = \frac{g}{l} \left( 2 \pm \sqrt{2} \right).
\end{equation}
These two frequencies, $\omega_+$ and $\omega_-$, represent the two fundamental modes of vibration for the double pendulum.

\subsection{Physical Implications of Normal Modes}
The normal modes represent specific patterns of motion where all parts of the system oscillate with the same frequency and a fixed phase relation. 

\begin{itemize}[leftmargin=*]
    \item \textbf{The In-Phase Mode (Lower Frequency, $\omega_-$):} Corresponding to $\omega_-^2 = (2 - \sqrt{2})\frac{g}{l}$, the amplitude ratio $A_2/A_1$ is positive (specifically, $\sqrt{2}$). Physically, this means both pendulum masses swing in the same direction simultaneously. The lower string remains relatively aligned with the upper string, reducing the overall restoring force from gravity. Consequently, this mode oscillates more slowly.
    \item \textbf{The Out-of-Phase Mode (Higher Frequency, $\omega_+$):} Corresponding to $\omega_+^2 = (2 + \sqrt{2})\frac{g}{l}$, the amplitude ratio $A_2/A_1$ is negative ($-\sqrt{2}$). In this mode, the two masses swing in opposite directions. This motion significantly bends the joint between the two rods, maximizing the restoring forces, which leads to a much faster, higher-frequency oscillation.
\end{itemize}
According to the superposition principle, any small-amplitude motion of the double pendulum can be described as a linear combination of these two fundamental normal modes. This theoretical foundation is crucial for verifying the accuracy of the numerical simulations presented in the following sections.

\subsection{Superposition State in the Linear Regime}
In the small-angle linear regime, the general solution is the superposition of the two normal modes:
\begin{equation}
\begin{aligned}
\theta_1(t) &= c_- A_{1,-}\cos(\omega_- t + \phi_-) + c_+ A_{1,+}\cos(\omega_+ t + \phi_+), \\
\theta_2(t) &= c_- A_{2,-}\cos(\omega_- t + \phi_-) + c_+ A_{2,+}\cos(\omega_+ t + \phi_+),
\end{aligned}
\end{equation}
where $c_{\pm}$ are mode amplitudes set by initial conditions, $\phi_{\pm}$ are initial phases, and $(A_{1,\pm}, A_{2,\pm})$ are eigenvectors associated with each normal mode. This form shows that the observed trajectory is not usually a pure in-phase or out-of-phase oscillation, but a mixed state containing both. When both coefficients are non-zero, energy is distributed across modes and the measured motion may show modulation and phase-mixing effects. The superposition description remains accurate only while linear assumptions hold; for larger amplitudes, non-linear coupling progressively invalidates pure modal superposition and leads to strongly irregular motion.

\subsection{Active Control and Rapid Amplitude Reduction}
In practical engineering applications, it is often desirable to suppress the oscillations of a double pendulum rapidly. Relying solely on passive natural damping (such as air drag and joint friction) is typically insufficient. By introducing an external control torque, we can actively extract energy from the system. Below are several theoretical approaches to achieve this in the small-angle regime.

\subsubsection{Active Damping via Velocity Feedback}
The most direct method is analogous to adding a viscous damper to a single pendulum. The linearized multi-degree-of-freedom equations of motion can be written in matrix form as:
\begin{equation}
M \ddot{\mathbf{\theta}} + C \dot{\mathbf{\theta}} + K \mathbf{\theta} = \mathbf{u}(t),
\end{equation}
where $M, C, K$ are the mass, damping, and stiffness matrices respectively, $\mathbf{\theta} = [\theta_1, \theta_2]^T$, and $\mathbf{u}$ is the control input. If we apply a control law $\mathbf{u} = -K_v \dot{\mathbf{\theta}}$ (where $K_v$ is a positive definite velocity gain matrix), we effectively increase the system's damping matrix to $C + K_v$. This active damping simultaneously dissipates energy from both normal modes, causing the vibrations to decay exponentially. While simple to implement, excessively high gains can lead to an overdamped, sluggish response.

\subsubsection{Anti-Phase Excitation and Dynamic Vibration Absorption}
Another approach draws inspiration from driven single pendulums. In a double pendulum, we can utilize one of the masses to suppress the vibration of the other. 
If an actuator is placed at the first pivot, it can apply a force perfectly out-of-phase with the measured velocity of the first pendulum, essentially acting as localized active damping.
A more elegant strategy is to use the second pendulum as an \emph{active dynamic vibration absorber} (similar to a Tuned Mass Damper, TMD). By controlling the torque at the second joint, the virtual natural frequency of the second pendulum can be tuned to match the current vibration frequency of the primary system. The vibrational energy of the first pendulum is rapidly transferred to the second pendulum, where it can be dissipated by the joint's damping.

\subsubsection{Optimal Control (LQR) and Modal Decoupling}
Because the double pendulum has two distinct normal modes (low-frequency in-phase and high-frequency out-of-phase), simple single-joint feedback might inadvertently suppress one mode while exciting the other. To address this modal coupling, modern control theory offers the Linear Quadratic Regulator (LQR).

Defining the state vector as $\mathbf{x} = [\theta_1, \theta_2, \dot{\theta}_1, \dot{\theta}_2]^T$, the system becomes $\mathbf{\dot{x}} = A\mathbf{x} + B\mathbf{u}$. The LQR method seeks a state-feedback control law $\mathbf{u} = -K\mathbf{x}$ that minimizes the cost function:
\begin{equation}
J = \int_{0}^{\infty} \left( \mathbf{x}^T Q \mathbf{x} + \mathbf{u}^T R \mathbf{u} \right) dt.
\end{equation}
By solving the Algebraic Riccati Equation, we obtain the optimal gain matrix $K$. This shifts the closed-loop poles deeper into the left half of the complex plane, ensuring the fastest possible decay of both modes simultaneously without exceeding the physical limits of the actuators. If both joints are equipped with motors and sensors, LQR is the most reliable and fastest-converging engineering solution for rapid amplitude reduction.

\section{Simulation}
\subsection{Using a Simple Pendulum to Identify Air Drag}
The model uses the \textbf{underdamped assumption}: $c^2<4mk$, with small-angle motion.
\begin{equation}
m\ell^2\ddot{\theta}+c\dot{\theta}+mg\ell\,\theta=0
\end{equation}
where $\ell$ is pendulum length, $m$ is bob mass, $c$ is air-drag coefficient, $g$ is gravitational acceleration, $\theta$ is angular displacement, and $t$ is time.

Standard form:
\begin{equation}
\ddot{\theta}+2a\dot{\theta}+w_0^2\theta=0,\qquad
a=\frac{c}{2m\ell^2},\qquad
w_0=\sqrt{\frac{g}{\ell}}
\end{equation}
where $a$ is damping parameter and $w_0$ is undamped angular frequency.

Under the underdamped condition, the solution is
\begin{equation}
\theta(t)=A e^{-at}\cos(wt+\phi),\qquad
w=\sqrt{w_0^2-a^2}
\end{equation}
where $A$ is amplitude constant, $\phi$ is initial phase, and $w$ is damped angular frequency.

With period-only measurement, use the damped period $T$:
\begin{equation}
T=\frac{2\pi}{w},\qquad
T_0=2\pi\sqrt{\frac{\ell}{g}}
\end{equation}
where $T$ is measured damped period and $T_0$ is theoretical undamped period.

Infer damping from period:
\begin{equation}
a=\sqrt{\left(\frac{2\pi}{T_0}\right)^2-\left(\frac{2\pi}{T}\right)^2}
\end{equation}
Then compute air-drag coefficient:
\begin{equation}
c=2m\ell^2 a
\end{equation}

\subsection{What to Measure and How to Measure}
\begin{equation}
\ell,\ m,\ N,\ t_{0,i},\ t_{N,i}\quad (i=1,\ldots,M)
\end{equation}
where $N$ is number of full periods per timing run, $M$ is number of trials, $t_{0,i}$ is start time of trial $i$, and $t_{N,i}$ is the time at the end of the $N$th period in trial $i$.

Period in each trial:
\begin{equation}
T_i=\frac{t_{N,i}-t_{0,i}}{N}
\end{equation}
Mean period:
\begin{equation}
\bar{T}=\frac{1}{M}\sum_{i=1}^{M}T_i
\end{equation}
Standard deviation (repeatability check):
\begin{equation}
s_T=\sqrt{\frac{1}{M-1}\sum_{i=1}^{M}(T_i-\bar{T})^2}
\end{equation}

Final estimates:
\begin{equation}
\bar{a}=\sqrt{\left(\frac{2\pi}{T_0}\right)^2-\left(\frac{2\pi}{\bar{T}}\right)^2},
\qquad
\bar{c}=2m\ell^2\bar{a}
\end{equation}
Use $\bar{c}$ as the damping coefficient in the double-pendulum simulation.

\subsection{Experimental Process}
\begin{equation}
\text{Step 1: measure }\ell,m
\rightarrow
\text{Step 2: set small angle }(\theta_0\le 10^\circ)
\rightarrow
\text{Step 3: release without push}
\end{equation}
\begin{equation}
\text{Step 4: record }t_{0,i},t_{N,i}\text{ for }M\text{ trials}
\rightarrow
\text{Step 5: compute }T_i,\bar{T},s_T
\rightarrow
\text{Step 6: compute }\bar{a},\bar{c}
\end{equation}
\section{Experiment Results}
  Experiment results content...

  \section{Conclusion}
  Conclusion content...

% References (AQA Harvard Style)
\section{References}
\begin{thebibliography}{99}
  % Book Reference
  \bibitem{book1}
  Smith, J. (2020). \emph{Title of the Book}. Publisher Name.
  
  % Journal Article Reference
  \bibitem{journal1}
  Johnson, A. B. (2019). Title of the article. \emph{Journal Name}, 15(3), 45-67.
  
  % Website Reference
  \bibitem{web1}
  Organisation Name. (2021). Title of web page. Retrieved from https://www.example.com
  
  % Chapter in Edited Book
  \bibitem{chapter1}
  Brown, C. (2018). Title of chapter. In A. Editor \& B. Co-editor (Eds.), \emph{Title of Book} (pp. 123-145). Publisher Name.
  
  % Conference Paper
  \bibitem{conference1}
  Wilson, D. (2017). Title of paper. In \emph{Proceedings of the 10th Conference on Topic} (pp. 67-78). Conference Location.
\end{thebibliography}

% Appendices (Optional)
\appendix
\section{Appendix A}
  Appendix content...

\end{document}
